../cictikz/src/cictikz/data/tex/cictikz_preamble.tex

% Current starved ring oscillator: a PMOS mirror copies I_control into
% the supply of each inverter, and a V_control PMOS in series with each
% branch throttles it further. One column per stage, no crossings.

\begin{circuitikz}[american, thick, transform shape, circuit ee IEC]

  ../cictikz/src/cictikz/data/tex/cictikz_lib.tex

  \def\xa{2.6}    % stage columns
  \def\xb{5.2}
  \def\xc{7.8}
  \def\ymir{5.6}  % mirror row (PMOS drain at \ymir, source at +\grid)
  \def\yctl{4.0}  % V_control row
  \def\yinv{1.4}  % the ring
  \def\ycap{0.0}

  % ---- the mirror: diode device fed by I_control ----
  \draw (0,{\ymir+\grid}) \vsupply;
  \draw (0,\ymir) \lvmpmos{MD}{};
  \coordinate (gmd) at (MD.G);
  \fill (gmd) circle (0.075);
  \draw (gmd) |- (0,\ymir);
  \fill (0,\ymir) circle (0.075);
  \draw (0,\ymir) to [I, l_=$I_{control}$] (0,{\ymir-2.0});
  \draw (0,{\ymir-2.0}) \vground;

  % mirror outputs, one per stage
  \draw (\xa,{\ymir+\grid}) \vsupply;
  \draw (\xa,\ymir) \lvpmos{M1}{};
  \draw (\xb,{\ymir+\grid}) \vsupply;
  \draw (\xb,\ymir) \lvpmos{M2}{};
  \draw (\xc,{\ymir+\grid}) \vsupply;
  \draw (\xc,\ymir) \lvpmos{M3}{};

  % common gate rail
  \coordinate (g1) at (M1.G);
  \coordinate (g2) at (M2.G);
  \coordinate (g3) at (M3.G);
  \draw (gmd) -- (g1);
  \draw (g1) -- (g2);
  \draw (g2) -- (g3);
  \fill (g1) circle (0.075);
  \fill (g2) circle (0.075);

  % ---- V_control PMOS in series with each branch ----
  \draw (\xa,\yctl) \lvmpmos{M4}{};
  \draw (\xb,\yctl) \lvmpmos{M5}{};
  \draw (\xc,\yctl) \lvmpmos{M6}{};
  \coordinate (g4) at (M4.G);
  \coordinate (g5) at (M5.G);
  \coordinate (g6) at (M6.G);
  \draw (g4) -- (g5);
  \draw (g5) -- (g6);
  \fill (g5) circle (0.075);
  \fill (g6) circle (0.075);
  \draw (g6) -- ++(0.9,0) coordinate (vctap);
  \draw (vctap) to [short,-o] ++(0.4,0) node[anchor=west] {$V_{control}$};

  % branch wiring: mirror drain to control source
  \draw (\xa,\ymir) -- (\xa,{\yctl+\grid});
  \draw (\xb,\ymir) -- (\xb,{\yctl+\grid});
  \draw (\xc,\ymir) -- (\xc,{\yctl+\grid});

  % ---- the ring ----
  \draw ({\xa-0.9},\yinv) node[not port] (i1) {};
  \draw ({\xb-0.9},\yinv) node[not port] (i2) {};
  \draw ({\xc-0.9},\yinv) node[not port] (i3) {};

  % each inverter's supply comes from its starved branch
  \draw (\xa,\yctl) -- (\xa,{\yinv+1.1}) -- ({\xa-0.9},{\yinv+1.1}) -- (i1.north);
  \draw (\xb,\yctl) -- (\xb,{\yinv+1.1}) -- ({\xb-0.9},{\yinv+1.1}) -- (i2.north);
  \draw (\xc,\yctl) -- (\xc,{\yinv+1.1}) -- ({\xc-0.9},{\yinv+1.1}) -- (i3.north);

  % stage to stage, caps on each output
  \draw (i1.out) -- (i2.in 1);
  \draw (i2.out) -- (i3.in 1);
  \coordinate (o1) at ({\xa+0.75},\yinv);
  \coordinate (o2) at ({\xb+0.75},\yinv);
  \fill (o1) circle (0.075);
  \fill (o2) circle (0.075);
  \draw (o1) to [C, l=$C$] ({\xa+0.75},{\ycap-0.9});
  \draw ({\xa+0.75},{\ycap-0.9}) \vground;
  \draw (o2) to [C, l=$C$] ({\xb+0.75},{\ycap-0.9});
  \draw ({\xb+0.75},{\ycap-0.9}) \vground;
  \coordinate (o3) at ({\xc+0.65},\yinv);
  \fill (o3) circle (0.075);
  \draw (o3) to [C, l_=$C$] ({\xc+0.65},{\ycap-0.9});
  \draw ({\xc+0.65},{\ycap-0.9}) \vground;

  % feedback
  \draw (o3) -- ({\xc+1.4},\yinv) -- ({\xc+1.4},{\ycap-1.7})
        -- ({\xa-1.9},{\ycap-1.7}) -- ({\xa-1.9},\yinv) -- (i1.in 1);

\end{circuitikz}
\end{document}
